\documentclass[10pt]{article}
\usepackage[margin=0.75in]{geometry}
\usepackage{amsmath,microtype}
\title{Guide: Replacing an Orbit Prefix}
\author{Collatz proof project}
\date{September 5, 2026}
\begin{document}
\maketitle
\raggedright
\section*{Why merging helps}
Imagine the smallest positive integer with an unbounded Collatz orbit.
If a smaller positive integer eventually joins that orbit, its orbit is
also unbounded, contradicting the choice of the smallest one. This is a
valid proof strategy only if we can guarantee that smaller predecessor.
Lean now proves the implication; existence remains open.

\section*{What the replacement certificate checks}
A finite odd/even pattern determines an affine endpoint formula. Two
patterns with the same length and number of odd steps have the same
multiplier. Increasing the correction by $3^j\delta$ reduces the starting
seed by $\delta$ while keeping the endpoint fixed. The proof checks that
the new integer actually follows every specified parity. An arbitrary
number in a plausible correction interval is insufficient.

The replacement is useful when $0<\delta<n$, so its seed is positive
and smaller. It need not grow at every intermediate step.

\section*{A working infinite family}
Any initial run of odd steps followed by two even steps can be replaced
by a pattern starting from the integer immediately below the original
seed. For example,
\[
\begin{array}{rrrrrrr}
15&23&35&53&80&40&20\\
14&7&11&17&26&13&20
\end{array}
\]
These are six shortcut steps with four odd steps on each row. This is
the classical Garner stem family, now formalized in this project.
It is not a new mathematical discovery or a universal coverage result.

\section*{Why this version cannot cover every seed}
Lean proves that 27 never merges with a smaller positive seed at the
same time with the same number of odd steps. This holds for every time,
not just a search cutoff. The proof checks the initial 70 steps and then
uses the repeating cycle: 27 has accumulated too many odd steps for any
smaller seed to catch up under these two restrictions.

Since 27 reaches one, this does not disprove the proposed strategy for
a hypothetical least unbounded seed. It does mean that an unrestricted
claim of equal-time, equal-count replacement for every seed would be false.

\section*{What the experiment establishes}
At depth 18, exact enumeration finds replacements for 193,309 of all
262,144 binary cylinders. Among the 7,495 patterns whose every prefix
coefficient is at least one, it covers 1,391. The uncovered patterns are
not divergent orbits. The counts do not prove asymptotic or universal coverage.

\section*{Reproduce and continue}
From \texttt{lean/}, run \texttt{lake build Collatz.WordSurgery}.
From the root, run \texttt{python3 analysis/word\_surgery.py} and
\texttt{python3 -m unittest discover -s analysis -p test\_word\_surgery.py}.
Different word lengths or odd counts may help, but neither nondivergence
nor exclusion of all nontrivial cycles has been proved here.
\end{document}
